\documentclass[10pt,a4paper]{amsart}
\usepackage[margin=19mm]{geometry}
\usepackage{amsmath,amssymb,mathtools}
\usepackage[scr=rsfs,cal=euler]{mathalfa}
\usepackage{xfrac}
\usepackage[unicode,hidelinks,bookmarks=false]{hyperref}
\newcommand{\ie}{i.e.\ }
\newcommand{\cf}{cf.\ }
\newcommand{\resp}{respectively\ }
\newcommand{\Subsection}{\subsection}
\providecommand{\nobreakdash}{\nobreak\hbox{-}}
\setlength{\emergencystretch}{2em}
\multlinegap0pt
\usepackage{bm}
\usepackage{bbm}
\usepackage{dsfont}
\usepackage{xspace}
\usepackage{cancel}
\usepackage{mathtools}
\usepackage{amsthm}
\makeatletter
\@ifundefined{theorem}{\newtheorem{theorem}{Theorem}}{}
\makeatother
\title[The Partition Pairing Theorems I]{The Partition Pairing Theorems I}
\author{George Andrews, Manosij Ghosh Dastidar}
\date{Source: arXiv:2608.22291v1 (CC BY 4.0)}
\begin{document}
\maketitle

\noindent\emph{Source and licence.} The Partition Pairing Theorems I. Version arXiv:2608.22291v1 (CC BY 4.0), distributed under the Creative Commons Attribution 4.0 International licence, as recorded in the arXiv licence metadata for this exact version and shown on the abstract page https://arxiv.org/abs/2608.22291v1. Full rights notice: licenses/arxiv-2608.22291-CC-BY-4.0.txt.\par\smallskip

\noindent\textit{Editorial scope.} Counted unit: the Theorem whose source label is \texttt{thm:square-switch} in the article (source0.tex lines 1448--1455, source label \texttt{thm:square-switch}), delivered together with the article's own complete original proof of it (source0.tex lines 1457--1529, one \texttt{proof} environment of 73 lines). No mathematics is written, repaired or re-derived: statement and proof are the article's own text line for line. Required context retained uncounted: thm:odd-overpartition (lines 1231--1236); eq:A2-M (lines 1375--1379). Every definition, display and label of those lines is retained unchanged. Omitted from this package and not redistributed: 1--1230, 1237--1374, 1380--1447, 1456--1456, 1530--1810. Nothing inside a retained excerpt is omitted or altered: every retained body is delivered exactly as the article's lines are written, so no omission receipt of any kind accompanies this package. Citations: the retained excerpts carry no citation command at all, so no bibliography entry of the article is transcribed and no reference is redistributed. Reference adaptation: none was needed and none was made; every reference inside the delivered unit resolves inside it, so no unresolved reference or citation is produced. Printed slips kept literal: no printed word, symbol, hypothesis or number of the counted statement or of its proof was altered. Layout and class: the article's own class is \texttt{article} with packages including geometry, amsmath, amssymb, amsthm, mathtools, microtype, enumitem, hyperref; the wrapper is the lane's pinned \texttt{amsart} editorial wrapper at 10pt on A4, which restates the theorem-like environments the retained text needs and adds the article's own macro definitions that the retained text needs (none), with extra wrapper packages (bm, bbm, dsfont, xspace, cancel, mathtools). The delivered numbering is the unit's own. Assets: no retained span contains a figure environment, a graphics-inclusion command, a PDF-page inclusion, a table or a TikZ picture; no image file is copied into this package, the package declares no asset dependency and no image file is redistributed with it. Licence: version arXiv:2608.22291v1 of this article is distributed under the Creative Commons Attribution 4.0 International licence, as recorded in the arXiv licence metadata for this exact version and shown on the abstract page https://arxiv.org/abs/2608.22291v1. A full rights notice is delivered at licenses/arxiv-2608.22291-CC-BY-4.0.txt. Characters: every retained excerpt is pure ASCII, so the delivered engine, XeLaTeX with its default Latin Modern text font, reports no missing character.
\begin{theorem}[Odd-overpartition theorem]\label{thm:odd-overpartition}
Let $M(n)$ be the number of simply paired partitions of $n$ having negative pairing rank.  Then, for every $n\geq1$,
\[
 M(n)=\frac12\,\overline p_{\mathrm{odd}}(n).
\]
\end{theorem}

\begin{equation}\label{eq:A2-M}
 (1+q^2)(1+q^4)A_2(q)
 =q(1-q^2)M(q)-\frac{q^3}{1+q}
 +\frac{q^6}{(1-q^2)^2}.
\end{equation}

\begin{theorem}[Square-switch theorem]\label{thm:square-switch}
For every $n\geq0$,
\begin{equation}\label{eq:square-switch}
 A_2(n+8)+A_2(n)
 \equiv n+1+\chi(n+3)+\chi(n+7)\pmod2.
\end{equation}
In particular, if $n$ is even, then $A_2(n+8)$ and $A_2(n)$ have the same parity if and only if one of $n+3$ and $n+7$ is a perfect square.
\end{theorem}

\begin{proof}
We first determine the parity of $M(n)$.  By Theorem~\ref{thm:odd-overpartition},
\begin{align*}
 1+2M(q)
 &=\prod_{j\geq1}
 \frac{1+q^{2j-1}}{1-q^{2j-1}}\\
 &=\prod_{j\geq1}
 \left(1+\frac{2q^{2j-1}}{1-q^{2j-1}}\right).
\end{align*}
Modulo $4$, every product containing at least two nonconstant factors vanishes.  Hence
\[
 1+2M(q)
 \equiv1+2\sum_{j\geq1}
 \frac{q^{2j-1}}{1-q^{2j-1}}\pmod4.
\]
Subtracting $1$ and dividing coefficientwise by $2$ gives
\begin{equation}\label{eq:M-parity}
 M(q)\equiv\sum_{j\geq1}
 \frac{q^{2j-1}}{1-q^{2j-1}}\pmod2.
\end{equation}
The coefficient of $q^m$ on the right of \eqref{eq:M-parity} is the number $d_{\mathrm{odd}}(m)$ of odd divisors of $m$.  If $m=2^\alpha v$ with $v$ odd, then
\[
 d_{\mathrm{odd}}(m)=d(v),
\]
which is odd precisely when $v$ is a square.  Therefore
\begin{equation}\label{eq:M-chi}
 M(m)\equiv\chi(m)\pmod2
 \qquad(m\geq0),
\end{equation}
where $M(0)=0$.

We define
\[
 \Delta_2(m)=A_2(m)+A_2(m-2)+A_2(m-4)+A_2(m-6).
\]
Modulo $2$, the two elementary series in \eqref{eq:A2-M} satisfy
\[
 -\frac{q^3}{1+q}+\frac{q^6}{(1-q^2)^2}
 \equiv\sum_{m\geq3}q^m+\sum_{j\geq0}q^{4j+6}\pmod2.
\]
We let $\eta(m)$ denote the coefficient of $q^m$ in the series on the right.  For $m\geq6$,
\begin{equation}\label{eq:eta}
 \eta(m)=
 \begin{cases}
 0,&m\equiv2\pmod4,\\
 1,&m\not\equiv2\pmod4.
 \end{cases}
\end{equation}
Taking coefficients in \eqref{eq:A2-M} and using \eqref{eq:M-chi}, we obtain
\begin{equation}\label{eq:Delta-parity}
 \Delta_2(m)
 \equiv\chi(m-1)+\chi(m-3)+\eta(m)\pmod2.
\end{equation}

The intermediate terms cancel in pairs modulo $2$, so
\begin{align*}
 A_2(n+8)+A_2(n)
 &\equiv\Delta_2(n+8)+\Delta_2(n+6)\\
 &\equiv\chi(n+7)+\chi(n+3)
   +\eta(n+8)+\eta(n+6)\pmod2.
\end{align*}
If $n$ is even, exactly one of $n+6$ and $n+8$ is congruent to $2$ modulo $4$, and \eqref{eq:eta} gives
\[
 \eta(n+8)+\eta(n+6)\equiv1\pmod2.
\]
If $n$ is odd, both $n+6$ and $n+8$ are odd, so both corresponding values of $\eta$ are $1$ and their sum is $0$ modulo $2$.  Thus, in all cases,
\[
 \eta(n+8)+\eta(n+6)\equiv n+1\pmod2.
\]
This proves \eqref{eq:square-switch}.

When $n$ is even, the integers $n+3$ and $n+7$ are odd, so $\chi$ is their ordinary perfect-square indicator.  Moreover, according to the residue of $n$ modulo $8$, at most one of these integers can be congruent to $1$ modulo $8$; hence they cannot both be squares.  It follows from \eqref{eq:square-switch} that the two values of $A_2$ have the same parity exactly when one of $n+3$ and $n+7$ is a square.
\end{proof}

\end{document}
